\section*{2. Lefschetz 线束: 结果的陈述}

\subsubsection*{2.1.}
以 $\mathbb P^r$ 表示 $k$ 上的 $r$ 维射影空间, $\check{\mathbb P}^{r}$ 表示由 $\mathbb P^r$ 中超平面构成的\underline{对偶}射影空间, $\operatorname{Gr}(1,\check{\mathbb P}^{r})$ 表示 $\check{\mathbb P}^{r}$ 中直线所成的簇. 点 $D\in\operatorname{Gr}(1,\check{\mathbb P}^{r})$ 称为 $\mathbb P^r$ 中的一个\underline{超平面束}(pencil of hyperplanes). 将 $D$ 视为 $\check{\mathbb P}^{r}$ 中的一条直线, 滥用记号写成 $D=\{H_t\}_{t\in D}$, 其中 $H_t$ 为作为直线 $D$ 的点的那些 $\mathbb P^r$ 中的超平面.

\underline{$D$ 的轴}为 $\mathbb P^r$ 中余维二的线性子簇, 即此束任意两个不同成员的交 $H_0\cap H_\infty$. 将此轴视作 $\operatorname{Gr}(r-2,\mathbb P^r)$ 的一点(此簇由 $\mathbb P^r$ 中的 $(r-2)$ 维线性子簇构成), 则把 $D$ 映到其轴的映射就是熟知的正交同构
\[
\operatorname{Gr}(1,\check{\mathbb P}^{r})
\xrightarrow{\ \sim\ }
\operatorname{Gr}(r-2,\mathbb P^r).
\]

\subsubsection*{2.2.}
设 $X$ 为维数 $n\geq1$ 的固有光滑不可约 $k$-概形, 并配有嵌入 $X\hookrightarrow\mathbb P^r$. $\mathbb P^r$ 中的超平面束 $D=\{H_t\}_{t\in D}$ 称为一个\underline{Lefschetz 线束}(相对于此嵌入), 如果:
\begin{enumerate}
\item[(a)] $D$ 的轴与 $X$ 横截相交(EGA IV, 17.13.7);
\item[(b)] 存在稠密开子集 $U\subset D$, 使对每个 $t\in U$, 超平面 $H_t$ 与 $X$ 横截相交;
\item[(c)] 对 $t_0\in D-U$, 超平面 $H_{t_0}$ 除一个点外与 $X$ 横截相交, 且该点为 $X\cap H_{t_0}$ 的普通二次奇点.
\end{enumerate}
下面将在 3.2.1 中看到, Lefschetz 线束构成 $\operatorname{Gr}(1,\check{\mathbb P}^{r})$ 的一个\underline{开子集}.

\subsubsection*{2.3.}
嵌入 $X\hookrightarrow\mathbb P^r$ 称为 \underline{Lefschetz 嵌入}, 如果 Lefschetz 线束构成 $\operatorname{Gr}(1,\check{\mathbb P}^{r})$ 的一个稠密开子集.
\footnote{更准确地说, 它们在 $\operatorname{Gr}(1,\check{\mathbb P}^{r})$ 的闭点集合中构成稠密开子集. 下面将允许这种语言上的滥用.}

\subsubsection*{2.4.}
回顾\underline{给定嵌入 $X\hookrightarrow\mathbb P^r$ 的倍数}这一概念. 它是该嵌入与第 $d$ 次 Veronese 嵌入
\footnote{印本称其为第 $d$ 次 Segre 嵌入; 所列的单项式映射为 Veronese 嵌入.}
\[
S_d:\mathbb P^r\hookrightarrow
\mathbb P^{\binom{r+d}{d}-1},
\qquad
(X_0,\ldots,X_r)\longmapsto(\ldots,X^W,\ldots),
\]
的复合, 其中
\[
X^W=X_0^{W_0}\cdots X_r^{W_r}
\quad\text{及}\quad
W_0+\cdots+W_r=d.
\]
因此, 对给定嵌入而言的“$X$ 与 $d$ 次超曲面的截面”---英语称为次数 $d$ 的 hypersurface section---就是此嵌入的第 $d$ 倍嵌入的一个超平面截面.

\subsubsection*{定理 2.5.}
\underline{设 $X$ 为 $k$ 上维数 $n\geq1$ 的不可约、固有、光滑概形, 并配有嵌入 $i:X\hookrightarrow\mathbb P^r$. 则:}
\begin{enumerate}
\item[2.5.1.] \underline{对 $d\geq2$, $i$ 的每个第 $d$ 倍嵌入均为 Lefschetz 嵌入;}
\item[2.5.2.] \underline{若 $k$ 的特征为零, 则给定嵌入 $i$ 是 Lefschetz 嵌入.}
\end{enumerate}
证明分成下面的几个部分.